\section{Business Process Modelling}
\begin{exercise}{1}
Using only the base BPMN elements (activities, gateways and sequence flows), model the following process:
\\\\An order-to-cash process is triggered by the receipt of a purchase order from a customer. 
Upon receipt, the purchase order has to be checked against the stock to determine if the requested item(s) are available.
Depending on stock availability thepurchase order may be confirmed or rejected.
If the purchase order is rejected, the process completes.
If the purchase order is confirmed, an invoice is emitted, and the goods requested are shipped.
\end{exercise}
\begin{comment}
    #TODO: Insert solution BPMN diagram here
    This was taken from lecture 2- half of document
\end{comment}

\begin{exercise}{2}
Using only the base BPMN elements (activities, gateways and sequence flows), model the following process:
\\\\Once a loan application has been approved, an
acceptance pack is prepared and sent to the customer.
The acceptance pack includes a repayment schedule,
which the customer needs to agree upon. The payment
documents are received back by the loan provider. The
latter then verifies the repayment agreement: if the
applicant disagreed with the repayment schedule, the
loan provider cancels the application; if the applicant
agreed, the loan provider approves the application. In
either case, the loan completes with the loan provider
notifying the applicant of the application status.
\end{exercise}

\begin{exercise}{3}
Given the following process description, create a BPMN model that captures all the details:
\\\\The clients of the IT helpdesk are employees of a company. A request may be an IT-
related problem that a client has, or an access request (e.g., requesting rights to access
an IT system). Requests need to be handled according to their type and their priority.
A client makes a request to the help desk. The help desk staffed with “Level-1” support
staff registers the request. Level-1 support staff typically are junior people with less than
12 months experience, but are capable of resolving simple requests. When the Level-1
employee knows how to resolve a request, she resolves it, notifies the resolution to the
client and closes the simple request.
When the Level-1 employee does not know how to resolve a request, the request is for-
warded to a more experienced “Level-2” support staff. When a Level-2 employee receives
a request, the request is evaluated, a resolution is developed and sent back to the Level-1
employee. Eventually, the Level-1 employee forwards the resolution to the client who
can test the resolution and notify the outcome to the IT helpdesk. The client notifies
the outcome of the test to the Level-1 employee via e-mail. If the client states that the
request is fixed, the request is closed successfully, and the process ends. If the request is
not fixed, it is resent to Level-2 support for further action and goes through the process
again. If, instead, after the client notification, no test outcome is received back for a
month, the Level-1 employee closes the request for timeout and the process ends.
\end{exercise}
\begin{comment}
    #TODO: Insert solution BPMN diagram here
    This was taken from exam mockup
\end{comment}

\begin{exercise}{4}
Given the following process description, create a BPMN model that captures all the details:
\\\\After an expense report is received from an employee, the report is then reviewed for the approval. 
If the amount is lower than 1000 euros, it is automatically approved. 
If the amount is equal or higher than 1000 euros, then a manual check is required. 
In case the report is not ok, it is rejected and the employee has to be notified via email. 
In case it is fine, it is approved, the reimbursement is deposited to the employee’s bank account and the
documentation properly archived in parallel.
\end{exercise}
\begin{comment}
    #TODO: Insert solution BPMN diagram here
    Taken from lesson 5 hands-on session
\end{comment}

\begin{exercise}{5}
Given the following process description, create a BPMN model that captures all the details:
\\\\When a new credit request from the customer is received by the credit assessor, the risk is assessed. 
If the risk is above a threshold, an advanced risk assessment has to be carried out,
otherwise a simple risk assessment will suffice. 
Once the assessment has been completed, the customer is notified with the result of the assessment
and meantime, the disbursement is organized.
For simplicity we assume that the result of an assessment is always positive.
\end{exercise}
\begin{comment}
    #TODO: Insert solution BPMN diagram here
    Taken from lesson 6 petri nets
\end{comment}
